\documentclass[../../../include/open-logic-section]{subfiles}

\begin{document}

\olfileid{his}{set}{hilbertcurve}

\olsection{પરિશિષ્ટ: હિલ્બર્ટની અવકાશ ભરતી વક્રરેખાઓ}

\olref[pathology]{sec} પ્રકરણમાં
આપણે જોયું કે રેખા અને
ચોરસમાં બિંદુઓની સંખ્યા
બરાબર છે એવી કૅન્ટૉરની
સાબિતીએ
(\olref[his][set][cantorplane]{thm:cantorplane})
પીઆનોને પૂછવા પ્રેર્યા કે
શું અવકાશ ભરતી
\emph{વક્રરેખા} હોઈ શકે.
1890માં તેમને તેનો
હકારાત્મક ઉત્તર મળ્યો.
આ વિભાગમાં હિલ્બર્ટની
અવકાશ ભરતી વક્રરેખા
કેવી રીતે રચવી તે
અનૌપચારિક રીતે
સમજાવીએ છીએ
(એક નાનકડા ફેરફાર સાથે).
\footnote{વધુ ચુસ્ત સમજૂતી
માટે \cite{Rose2010} જુઓ.
ફેરફાર એટલે નીચેની
વક્રરેખાઓના લાલ
ભાગો ઉમેરવા; તેથી
વક્રરેખા સતત છે તે
તપાસવું થોડું સહેલું બને છે.}

આપણે વિધેયો $h_1$, $h_2$, $h_3$,
\dots\@ના અનુક્રમના
લક્ષ તરીકે વિધેય $h$
વ્યાખ્યાયિત કરવાનું છે.
પહેલાં રચના વર્ણવીશું,
પછી તે અવકાશ ભરે છે
તે અને અંતે તે
વક્રરેખા છે તે બતાવીશું.

$h$નો વિસ્તાર એકમ ચોરસ
$\unitsquare$ લઈશું.
$h$નો પહેલો અંદાજ,
એટલે $h_1$, આ છે:
\begin{center}
	\begin{tikzpicture}
	\draw[gray] (0pt, 0pt) rectangle (80pt, 80pt);
	\draw[step = 40pt, gray, very thin] (0pt, 0pt) grid (80pt, 80pt);
	\draw[oldiagcolorC, thick] (20pt, 0)--(20pt, 20pt);
	\draw[oldiagcolorC, thick] (60pt, 0)--(60pt, 20pt);		
	\begin{scope}[xshift=20pt, yshift=20pt]
	\draw [black, thick, l-system={Hilbert curve, step=40pt, angle=90, axiom=L, order=1}]  lindenmayer system;
	\end{scope}
	\end{tikzpicture}
\end{center}
રચનાને અનુસરવા ચોરસ પર
$2 \times 2$ની જાળી મૂકી છે.
વક્રરેખા નીચેના ડાબા ચતુર્થાંશથી
શરૂ થઈ ઉપર ડાબે,
પછી ઉપર જમણે અને
અંતે નીચે જમણે જાય
છે એમ વિચારી શકીએ.
રચનાનો બીજો તબક્કો,
એટલે $h_2$, આ છે:
\begin{center}
	\begin{tikzpicture}
	\draw[gray] (0pt, 0pt) rectangle (80pt, 80pt);
	\draw[step = 20pt, gray, very thin] (0pt, 0pt) grid (80pt, 80pt);
	\draw[oldiagcolorC, thick] (0pt, 10pt)--(10pt, 10pt);
	\draw[oldiagcolorC, thick] (70pt, 10pt)--(80pt, 10pt);
	\draw[thick, oldiagcolorE] (10pt, 30pt)--(10pt, 50pt); 
	\draw[thick, oldiagcolorE] (30pt, 50pt)--(50pt, 50pt); 
	\draw[thick, oldiagcolorE] (70pt, 50pt)--(70pt, 30pt); \begin{scope}[xshift=10pt, yshift=30pt]
	\draw [thick, rotate=270,black, l-system={Hilbert curve, step=20pt, angle=90, axiom=L, order=1}]  lindenmayer system;
	\end{scope}
	\begin{scope}[xshift=10pt, yshift=50pt]
	\draw [thick, black, l-system={Hilbert curve, step=20pt, angle=90, axiom=L, order=1}]  lindenmayer system;
	\end{scope}
	\begin{scope}[xshift=50pt, yshift=50pt]
	\draw [thick, black, l-system={Hilbert curve, step=20pt, angle=90, axiom=L, order=1}]  lindenmayer system;
	\end{scope}
	\begin{scope}[xshift=70pt, yshift=10pt]
	\draw [thick, rotate=90,black, l-system={Hilbert curve, step=20pt, angle=90, axiom=L, order=1}]  lindenmayer system;
	\end{scope}
	\end{tikzpicture}
\end{center}
જુદા રંગો $h_2$ની
રચના સમજવામાં મદદ કરશે.
પહેલાં $h_1$ના
!!{colorC} સિવાયના
ભાગની નાની નકલો
ચોરસના નીચે ડાબા,
ઉપર ડાબા, ઉપર જમણા
અને નીચે જમણા
ચતુર્થાંશમાં મૂકીએ
(કાળા રંગમાં).
પછી આ ચાર આકૃતિઓને
!!{colorE} રેખાઓથી
જોડીએ. અંતે આકૃતિને
!!{colorC} રેખાઓથી
ચોરસની સીમા સાથે જોડીએ.

હવે $h_3$ તરફ વળીએ.
$h_1$ના લાલ સિવાયના
ભાગની જોડેલી ચાર નાની
નકલોથી $h_2$ બન્યું,
તે જ રીતે $h_2$ના
લાલ સિવાયના ભાગની
ચાર નાની નકલોથી
$h_3$ બને છે
(કાળા રંગમાં).
પછી તેમને !!{colorE}
રેખાઓથી એકબીજા
સાથે અને અંતે
!!{colorC} રેખાઓથી
ચોરસની સીમા
સાથે જોડીએ.
\begin{center}
	\begin{tikzpicture}
	\draw[gray] (0pt, 0pt) rectangle (80pt, 80pt);
	\draw[step = 10pt, gray, very thin] (0pt, 0pt) grid (80pt, 80pt);
	\draw[thick, oldiagcolorC](5pt, 0pt)--(5pt, 5pt); 
	\draw[thick, oldiagcolorC] (75pt, 0pt)--(75pt, 5pt); 
	\draw[thick, oldiagcolorE] (75pt, 45pt)--(75pt, 35pt);
	\draw[thick, oldiagcolorE] (5pt, 35pt)--(5pt, 45pt); 
	\draw[thick, oldiagcolorE] (35pt, 45pt)--(45pt, 45pt); 
	\draw[thick, oldiagcolorE] (75pt, 45pt)--(75pt, 35pt); \begin{scope}[xshift=5pt, yshift=35pt]
	\draw [rotate=270, thick, black, l-system={Hilbert curve, step=10pt, angle=90, axiom=L, order=2}]  lindenmayer system;
	\end{scope}
	\begin{scope}[xshift=5pt, yshift=45pt]
	\draw [thick, black, l-system={Hilbert curve, step=10pt, angle=90, axiom=L, order=2}]  lindenmayer system;
	\end{scope}
	\begin{scope}[xshift=45pt, yshift=45pt]
	\draw [thick, black, l-system={Hilbert curve, step=10pt, angle=90, axiom=L, order=2}]  lindenmayer system;
	\end{scope}
	\begin{scope}[xshift=75pt, yshift=5pt]
	\draw [thick, rotate=90,black, l-system={Hilbert curve, step=10pt, angle=90, axiom=L, order=2}]  lindenmayer system;
	\end{scope}
	\end{tikzpicture}
\end{center}
આ રીતે $h_n$ પરથી
$h_{n+1}$ વ્યાખ્યાયિત
કરવાની સામાન્ય રીત
દેખાય છે. હવે આ
અનુક્રમી વિધેયો
$h_1$, $h_2$, \dots\@
ના બિંદુવાર લક્ષથી
વક્રરેખા $h$ને
\emph{પોતે} વ્યાખ્યાયિત
કરીએ. એટલે દરેક
$x \in \unitline$ માટે:
\begin{align*}
	h(x) &= \lim_{n \rightarrow \infty} h_n(x)
\end{align*} 
\sourcecorrection{OLMD-170}{મૂળમાં
$x \in \unitsquare$ લખાયું છે,
પણ $h_n$ અને $h$નો
પ્રદેશ $\unitline$ હોવો
જોઈએ. ચિત્રો માત્ર
વક્રરેખાઓના માર્ગ
દર્શાવે છે; દરેક
$h_n$નું પરિમાણીકરણ
અને તેનું બિંદુવાર
લક્ષ અસ્તિત્વમાં છે
તે અહીં સાબિત
થતું નથી.}
હવે આ વક્રરેખા
અવકાશ ભરે છે
તે બતાવીએ.
$h_n$ દોરીએ ત્યારે
$\unitsquare$ પર
$2^n \times 2^n$ની
જાળી મૂકીએ છીએ.
પાયથાગોરસના પ્રમેયથી
જાળીના દરેક કોષના
કર્ણની લંબાઈ છે:
\[
\sqrt{\left(\nicefrac{1}{2^{n}}\right)^2+\left(\nicefrac{1}{2^{n}}\right)^2} = 2^{(\frac{1}{2}-n)}
\]
અને સ્પષ્ટપણે $h_n$
જાળીના દરેક કોષમાંથી
પસાર થાય છે. તેથી
$\unitsquare$નો દરેક
બિંદુ $h_n$ પરના
કોઈ બિંદુથી \emph{વધુમાં વધુ}
$2^{(\frac{1}{2}-n)}$
અંતરે છે. હવે $h$ને
વિધેયો $h_1$, $h_2$,
$h_3$, \dots\@ના લક્ષ
તરીકે વ્યાખ્યાયિત કર્યું છે.
તેથી મૂળની દલીલ પ્રમાણે
કોઈ પણ બિંદુનું $h$થી
મહત્તમ અંતર આ લક્ષ છે:
\[
\lim_{n \rightarrow \infty} 2^{(\frac{1}{2}-n)} = 0.
\]
અર્થાત્, $\unitsquare$નો
દરેક બિંદુ $h$થી
$0$ અંતરે છે. બીજા
શબ્દોમાં, $\unitsquare$નો
દરેક બિંદુ વક્રરેખા
\emph{પર} છે; તેથી
$h$ અવકાશ ભરે છે!{}
\sourcecorrection{OLMD-171}{દરેક
$h_n$ની છબી જાળીના
બધા કોષને મળે છે અને
માત્ર બિંદુવાર લક્ષ
છે, તે પરથી છબી
$h(\unitline)$ આખો
ચોરસ ભરે છે એવો
નિષ્કર્ષ નીકળતો નથી.
તે માટે વિધેયોનું
સમરૂપ અભિસરણ
સ્થાપવું, અને
લક્ષની છબી બંધ છે
તેવું સાતત્ય/સઘનતા
આધારિત પગલું
આપવું પડે.
આ બંને અહીં
અસાબિત છે.}

હવે $h$ ખરેખર
\emph{વક્રરેખા} છે
તે બતાવવાનું બાકી છે.
એ માટે આ ખ્યાલ
વ્યાખ્યાયિત કરવો પડે.
આધુનિક વ્યાખ્યા
જોર્ડને 1887માં
આપેલી વ્યાખ્યા
પર આધારિત છે
(પહેલી અવકાશ ભરતી
વક્રરેખા રજૂ થાય
તેના થોડાં વર્ષ પહેલાં):

\begin{defn}
વક્રરેખા એટલે
$\unitline$માંથી
$\Real^2$માંનું
સતત ચિત્રણ.
\end{defn}

આ વ્યાખ્યા ઘણી સહજ
લાગે છે: સહજ રીતે,
વક્રરેખા પ્રમાણભૂત
રેખાને સમતલ
$\Real^2$ પર લઈ
જતું “મસૃણ” ચિત્રણ
છે. આપણું વિધેય
$h$ ખરેખર
$\unitline$માંથી
$\Real^2$માં જાય છે.
તેથી $h$ સતત
છે તે બતાવવું બાકી છે.
\olref[limits]{sec}માં
આપણે $\epsilon$/$\delta$
સંકેતથી સાતત્ય
વ્યાખ્યાયિત કર્યું હતું.
મૂળની બોલચાલની
વાત આ છે:
\emph{જો તમે $\unitsquare$માં
બિંદુ $p$ અને ઇચ્છિત
ચોકસાઈ $\epsilon$ આપો,
તો $\unitline$નો એવો
ખુલ્લો ભાગ મળે કે
તેમાંના કોઈ પણ $x$
માટે $h(x)$ એ $p$થી
$\epsilon$ કરતાં ઓછા
અંતરે હોય.}
\sourcecorrection{OLMD-172}{અહીં
“મસૃણ” ગણિતીય
વ્યાખ્યા નથી; વ્યાખ્યા
માત્ર સાતત્ય માગે છે.
મૂળનું પછીનું બોલચાલનું
વિધાન પણ સાતત્ય
નથી બતાવતું: તે
દરેક લક્ષ્યબિંદુની
પાસે ક્યાંક આગતોનો
ખુલ્લો ભાગ મળે
એવું કહે છે. સાચી
સાતત્ય-શરત દરેક
$t\in\unitline$
અને $\epsilon>0$
માટે એવો $\delta>0$
માગે છે કે
$|x-t|<\delta$
હોય ત્યારે
$\|h(x)-h(t)\|<\epsilon$
હોય.}

હવે માનો કે તમે
$p$ અને $\epsilon$
આપ્યાં છે. તેનો
અર્થ $\unitsquare$માં
કેન્દ્ર $p$ અને
ત્રિજ્યા $\epsilon$
વાળું વર્તુળ દોરવા
જેવો છે. (વર્તુળનો
ભાગ $\unitsquare$ની
સીમા બહાર જઈ શકે,
પણ તેથી વાંધો નથી.)
યાદ કરો કે $h_n$
વર્ણવતાં આપણે
$\unitsquare$ પર
$2^n \times 2^n$ની
જાળી મૂકી હતી.
$\epsilon$ ગમે તેટલું
નાનું હોય, એવું
કોઈ $n$ મળે છે કે
$\unitsquare$ પરની આ
$2^n \times 2^n$ જાળીના ઓછામાં
ઓછો એક કોષ
કેન્દ્ર $p$ અને
ત્રિજ્યા $\epsilon$
વાળા વર્તુળની
અંદર આખો આવે છે.

એવું $n$ લો અને
$I$ને $\unitline$નો
એવો સૌથી મોટો
ખુલ્લો ભાગ માનો
જેને $h_n$ સંબંધિત
જાળીકોષની અંદર
આખો ચિતરે છે.
($2^n\times 2^n$
જાળીના દરેક કોષમાંથી
$h_n$ પસાર થાય
છે એમ આપણે
નોંધ્યું હોવાથી
એવો $(a,b)$
મળે છે.) હવે
જ્યારે પણ $x \in I$
હોય ત્યારે $h(x)$
એ જ જાળીકોષમાં
રહે છે તે બતાવવું
પૂરતું છે. અને
\emph{આ} માટે
દરેક $m > n$
માટે $h_m(x)$
એ જ કોષમાં
રહે છે તે
બતાવવું પૂરતું છે.
મૂળની દલીલમાં
આ સ્પષ્ટ ગણાય છે:
$m > n$ માટે
$h_m$માં એકમ
અંતરાલનો બરાબર
“એ જ ભાગ”
એ જ જાળીકોષમાં
ચિતરાય છે;
ફક્ત ત્યાં તેની
રચના વધુ લંબાયેલી
અને વળાંકવાળી
થતી જાય છે.
\sourcecorrection{OLMD-173}{આ
દલીલ આપેલા
$p$ની પાસે
વક્રરેખાનો થોડો
ભાગ આવે છે
એવું બતાવવાનો
પ્રયત્ન કરે છે;
પણ દરેક નિશ્ચિત
આગત $t$ની પાસે
$h(x)$ એ $h(t)$ની
પાસે રહે છે,
તે બતાવતી નથી.
વળી “સૌથી મોટો”
ખુલ્લો ભાગ અને
બધાં પછીના
$h_m$ની કોષ-સ્થિરતા
પરિમાણીકરણ
વિના સ્થાપિત
નથી. તેથી અહીં
સાતત્યની સાબિતી
હજી અધૂરી છે.}

\end{document}
